\documentclass[a4paper,fleqn]{cas-sc}

\usepackage[utf8]{inputenc}
\usepackage[T1]{fontenc}

\usepackage[authoryear,longnamesfirst]{natbib}

\usepackage{graphicx}
\graphicspath{{img/}}
\usepackage{booktabs}
\usepackage{array}
\usepackage{listings}
\usepackage{xcolor}
\usepackage{amsmath}
\usepackage{amssymb}
\usepackage{xurl}
\usepackage{multirow}
\usepackage{hyperref}
\usepackage{subfig}
\usepackage{float}
\floatplacement{figure}{!htbp}
\usepackage[section]{placeins}
\usepackage{capt-of}
\captionsetup{hypcap=false}
\usepackage{ragged2e}
\usepackage{needspace}

% --- Float placement tuning (keep figures near where they are referenced) ---
\setcounter{topnumber}{2}
\setcounter{bottomnumber}{1}
\setcounter{totalnumber}{3}
\renewcommand{\topfraction}{0.9}
\renewcommand{\bottomfraction}{0.8}
\renewcommand{\textfraction}{0.1}
\renewcommand{\floatpagefraction}{0.8}

\DeclareMathOperator{\Corr}{Corr}
\newcommand{\dCor}{\mathrm{dCor}}
\newcommand{\dCov}{\mathrm{dCov}}

\definecolor{codegreen}{rgb}{0,0.6,0}
\definecolor{codegray}{rgb}{0.5,0.5,0.5}
\definecolor{codepurple}{rgb}{0.58,0,0.82}
\definecolor{backcolour}{rgb}{0.95,0.95,0.92}

\lstdefinestyle{mystyle}{
  backgroundcolor=\color{backcolour},
  commentstyle=\color{codegreen},
  keywordstyle=\color{magenta},
  numberstyle=\tiny\color{codegray},
  stringstyle=\color{codepurple},
  basicstyle=\ttfamily\footnotesize,
  breakatwhitespace=false,
  breaklines=true,
  captionpos=b,
  keepspaces=true,
  numbers=left,
  numbersep=5pt,
  showspaces=false,
  showstringspaces=false,
  showtabs=false,
  tabsize=2
}
\lstset{style=mystyle}

\def\tsc#1{\csdef{#1}{\textsc{\lowercase{#1}}\xspace}}
\tsc{WGM}
\tsc{QE}
\tsc{EP}
\tsc{PMS}
\tsc{BEC}
\tsc{DE}

\begin{document}
\sloppy\setlength{\emergencystretch}{2em}
\let\WriteBookmarks\relax
\def\floatpagepagefraction{.8}
\def\textpagefraction{.1}
\justifying

\shorttitle{acedcor: Detecting Nonlinear Relationships}
\shortauthors{S. M. Park and H.-M. Kim}

\title [mode = title]{acedcor: A Python package for detecting and diagnosing nonlinear relationships using ACE and Distance Correlation}

\tnotetext[1]{This software paper describes the implementation, architecture, and core functionality of \texttt{acedcor}, a Python package designed to detect and diagnose complex nonlinear dependencies, built upon the ACE and Distance Correlation methodologies.}

\author[1]{{Sang Min Park}}[
  role=Conceptualization,
  bioid=1,
  orcid=0000-0000-0000-0000
]
\credit{Conceptualization of this study, Methodology, Software, Writing - Original draft preparation}

\affiliation[1]{organization={Department of Data Science, Konkuk University},
  addressline={120 Neungdong-ro, Gwangjin-gu},
  city={Seoul},
  postcode={05029},
  country={Republic of Korea}}

\author[2]{{Jun Zhao}}[
  role=Methodology,
  bioid=2,
  orcid=0000-0000-0000-0000
]
\ead{zhaojun2021@hotmail.com}
\credit{Formal analysis, Investigation, Validation}

\affiliation[2]{organization={School of Mathematics and Statistics, Ningbo University},
  addressline={818 Fenghua Road},
  city={Ningbo},
  postcode={315211},
  state={Zhejiang},
  country={China}}

\author[3]{{Hyoung-Moon Kim}}[
  role=Supervision,
  bioid=3,
  orcid=0000-0000-0000-0000
]
\cormark[1]
\ead{hmkim@konkuk.ac.kr}

\credit{Supervision, Funding acquisition, Project administration, Writing - Review \& Editing}

\affiliation[3]{organization={Department of Applied Statistics, Konkuk University},
  addressline={120 Neungdong-ro, Gwangjin-gu},
  city={Seoul},
  postcode={05029},
  country={Republic of Korea}}

\cortext[cor1]{Correspondence to: 120 Neungdong-ro, Gwangjin-gu, Seoul 05029, Republic of Korea.}


% [수정] 초록 (Park2025_ref 인용 제거 및 내용 명확화)
\begin{abstract}\sloppy
The Pearson correlation coefficient is a standard metric for linear dependence but often fails to capture complex nonlinear relationships inherent in real-world data. To address this limitation, we introduce \texttt{acedcor}, a Python package that integrates the Alternating Conditional Expectation (ACE) algorithm with Distance Correlation ($\dCor$). This unified framework optimally linearizes relationships before assessing dependence, significantly enhancing detection power. We demonstrate through extensive simulations that the ACE-$\dCor$ approach requires minimal assumptions and remains robust against heavy-tailed and skewed error distributions. Uniquely, the package provides a diagnostic metric, $\Delta\dCor$, which identifies hidden symmetric nonlinear structures (e.g., $y=x^2$) where traditional metrics fail. By making these advanced statistical methods accessible via a simple Python interface, \texttt{acedcor} empowers researchers across disciplines to accurately detect, diagnose, and model complex variable dependencies.
\end{abstract}

\begin{keywords}
distance correlation \sep nonlinear dependency \sep ACE algorithm \sep detection power \sep heavy-tailed distribution
\end{keywords}
\maketitle

\section*{Metadata}

\label{metadata}



\begin{table}[!htbp]
\small\sloppy\setlength{\emergencystretch}{2em}\centering
\caption{Code and Software Metadata Summary}
\label{tbl:unifiedMetadata}
\renewcommand{\arraystretch}{1.15}
\begin{tabular}{|>{\raggedright\arraybackslash}p{0.42\linewidth}|>{\raggedright\arraybackslash}p{0.48\linewidth}|}
\hline
\textbf{Description} & \textbf{Metadata} \\
\hline
\multicolumn{2}{|l|}{\textbf{MANDATORY CODE METADATA}} \\
\hline
Current code version & v1.0.0 \\
Permanent link to code/repository used for this code version & \url{https://github.com/ssangmin-junior/NonlinearCorr} \\
Code versioning system used & git \\
Software code languages, tools, and services used & Python, NumPy, Pandas, rpy2, R (acepack, energy), Matplotlib, SciPy, Scikit-learn \\
Compilation requirements, operating environments \& dependencies & Python (>=3.8), R (>=4.0), rpy2, NumPy, Pandas, energy, acepack, Matplotlib, SciPy, Scikit-learn \\
Support email for questions & sm123012@naver.com \\
\hline
\multicolumn{2}{|l|}{\textbf{OPTIONAL EXECUTABLE SOFTWARE METADATA}} \\
\hline
Current software version & v1.0.0 \\
Computing platforms/Operating Systems & Linux, OS X, Microsoft Windows \\
Installation requirements \& dependencies & (See repository link) \\
Support email for questions & sm123012@naver.com \\
\hline
\end{tabular}
\end{table}

\par\medskip
\normalsize
\raggedright
\justify

% ========== Two-Column Mode REMOVED (One-Column) ==========
% \twocolumn 명령어가 제거되었습니다.

\section{Motivation and significance}
\label{sec:motivation}
The primary task in data science is to identify and quantify the relationship between variables. The most common tool for this is the Pearson correlation coefficient ($\rho$), which is highly effective at measuring linear relationships \citep{CasellaBerger2002}. However, real-world data often exhibits complex nonlinear dependencies, which $\rho$ frequently fails to detect. For instance, a perfect $Y=X^{2}+Z$ parabolic relationship can yield a Pearson $\rho\approx0,$ leading to the false conclusion of "no relationship" \citep{CasellaBerger2002}. This inherent limitation in assessing general dependence necessitates the use of more robust measures.

To address this challenge, two key methodological streams have developed. The first is the use of robust dependence measures like Distance Correlation ($\dCor$) \citep{Szekely2007}, which is powerful because $\dCor(X,Y)=0$ if and only if $X$ and $Y$ are statistically independent, making it suitable for all relationship forms. The second is the use of non-parametric variable transformation techniques like the Alternating Conditional Expectation (ACE) algorithm \citep{BreimanFriedman1985}, which seeks to maximize the linear correlation between transformed variables, thereby estimating the true functional relationship.

The \texttt{acedcor} Python package integrates these two concepts into a novel ACE-$\dCor$ framework to provide a unified tool for detecting and diagnosing non-linear relationships. Specifically, the package implements the methodology that applies the non-parametric ACE transformation before calculating $\dCor$. This integrated approach not only confirms the existence of dependence but, through the diagnostic metric $\Delta\dCor$, provides crucial insight into the nature of the underlying nonlinear relationship, particularly symmetric patterns where traditional metrics fail.

\Needspace{16\baselineskip}
\begin{figure}
\centering
\subfloat[Before ACE transformation]{
  \includegraphics[width=0.49\textwidth]{Before_ace_plot.png}
}
\subfloat[After ACE transformation]{
  \includegraphics[width=0.49\textwidth]{after_ace_plot.png}
}
\caption{Example visualization before and after ACE transformation. Left: original relationship. Right: relationship after ACE-based transformation.}
\label{fig:intro_example}
\end{figure}

\section{Software description}
\label{sec:software_desc}

The \texttt{acedcor} package is a lightweight, open-source Python library built to seamlessly execute the ACE-$\dCor$ framework. It addresses the practical need for a unified tool that calculates, compares, and visualizes dependence measures both before and after optimal non-parametric transformation.

\subsection{Software Architecture}
The core architectural choice of \texttt{acedcor} is the reliance on the $\texttt{R}$ programming language for its essential computations, integrated seamlessly within a $\texttt{Python}$ environment via the $\texttt{rpy2}$ library. The codebase is structured for clarity and maintainability: \texttt{acedcor/core.py} contains the primary algorithms bridging Python and R, while \texttt{\_\_init\_\_.py} exposes the user-facing API. The \texttt{setup.py} handles dependency management, ensuring seamless integration of mandatory R packages. This decision prioritizes statistical rigor and computational correctness over a purely Python-native implementation.

The package's dependency stack consists of three layers. First, the Python Layer handles data management using $\texttt{NumPy}$ and $\texttt{Pandas}$, and visualization via $\texttt{Matplotlib}$. Second, the $\texttt{rpy2}$ Bridge serves as the crucial interface, allowing Python functions to call $\texttt{R}$ functions and pass data structures like $\texttt{NumPy}$ arrays back and forth efficiently. Finally, the $\texttt{R}$ Core Implementations provide the mandatory statistical backends: $\texttt{acepack}$ provides the seminal $\texttt{ace()}$ function for the Alternating Conditional Expectation algorithm \citep{BreimanFriedman1985}, and $\texttt{energy}$ provides the rigorously tested $\texttt{dcor()}$ function for Distance Correlation \citep{Szekely2007}.

Due to this architecture, the user must have a working $\texttt{R}$ installation, including the $\texttt{acepack}$ and $\texttt{energy}$ packages, on their system. This necessity ensures that \texttt{acedcor} leverages established, validated statistical backends for non-parametric optimal transformation and distance correlation, maintaining the integrity of the research methodology while offering accessibility through the dominant Python data science ecosystem.

\subsection{Software Functionalities}
To ensure accessibility, \texttt{acedcor} is designed to be easily installed via standard Python package managers. The package requires Python 3.8 or higher and an R environment (version 4.0 or higher) with the necessary libraries installed. Installation can be performed using pip commands as shown below.

\begin{lstlisting}[language=bash]
# 1. Install R packages (in R console)
# install.packages(c("acepack", "energy"))

# 2. Install acedcor
pip install acedcor
\end{lstlisting}

Listing \ref{list:basic_usage} demonstrates the basic workflow of the package. The function \texttt{calculate\_dcor\_improvement} takes two array-like inputs, $X$ and $Y$, and returns a dictionary containing the distance correlation before transformation, after ACE transformation, and the improvement metric ($\Delta\dCor$). This "Hello World" example allows users to immediately verify the package's operation using a simple parabolic dataset.

\begin{lstlisting}[language=Python, caption={Basic usage of acedcor to calculate $\Delta\dCor$.}, label={list:basic_usage}]
import numpy as np
from acedcor import calculate_dcor_improvement

# 1. Generate simple data (Parabola: y = x^2)
np.random.seed(42)
x = np.random.uniform(-1, 1, 500)
y = x**2 + np.random.normal(0, 0.1, 500)

# 2. Calculate Delta dCor
# This function handles ACE transformation internally
results = calculate_dcor_improvement(x, y)

print(f"dCor Before: {results['dcor_before']:.4f}")
print(f"dCor After:  {results['dcor_after']:.4f}")
print(f"Delta dCor:  {results['delta_dcor']:.4f}")

# Output Interpretation:
# A high Delta dCor (> 0.2) indicates a hidden symmetric structure.
\end{lstlisting}

The package is primarily centered around facilitating an integrated diagnostic workflow. The central component is the calculation of the Distance Correlation Improvement ($\Delta\dCor$ = $\dCor_{\text{After}} - \dCor_{\text{Before}}$). This metric serves as a quantitative diagnostic tool, where a significantly high $\Delta\dCor$ indicates a previously masked, structured relationship. The package is also equipped to handle comparative analyses with alternative, parametric transformations, such as the Box-Cox and Yeo-Johnson methods.

\Needspace{14\baselineskip}
\begin{center}
\captionof{table}{Input/Output schema for V3 diagnostic APIs.}
\label{tab:v3_io_schema}
\begin{tabular}{@{}p{0.26\textwidth}p{0.22\textwidth}p{0.44\textwidth}@{}}
\toprule
Function & Input & Output \\ \midrule
\texttt{dcor\_test} & \texttt{x, y, B, random\_state} & \texttt{(stat, pvalue)} \\
\texttt{acedcor\_test} & \texttt{x, y, B, random\_state} & \texttt{dcor\_before, p\_before, dcor\_after, p\_after, delta\_dcor} \\
\makecell[l]{\texttt{conditional\_}\\\texttt{dependence\_report}} & \texttt{x, y, z} & \makecell[l]{\texttt{dcor\_raw, dcor\_partial\_residual,}\\\texttt{partial\_distance\_correlation}} \\
\texttt{auto\_diagnosis\_report} & \texttt{x, y, output\_dir, B} & Summary dict + generated plots \\ \bottomrule
\end{tabular}
\end{center}

In addition, the automatic report function follows a fixed recommendation rule: when $\Delta\dCor > 0.2$, the report flags likely symmetric/non-monotonic structure and recommends ACE-based modeling; otherwise, it recommends simpler linear/monotonic transformations.


\section{Theoretical Background: Measures of Dependence in \texttt{acedcor}}
\label{sec:theory}
The \texttt{acedcor} package is built upon the rigorous combination of optimal transformation and robust dependence measurement. The following sections detail how these underlying statistical concepts are utilized within the package's functionality to effectively detect nonlinear relationships.

\subsection{Pearson Correlation Coefficient ($\rho$): The Baseline Limitation}
Pearson's $\rho$ is the most widely used measure of linear dependence. For a population of random variables $(X, Y)$, the coefficient is defined as:
$$\rho_{X,Y}=\frac{\mathrm{Cov}(X,Y)}{\sigma_{X}\sigma_{Y}}=\frac{E[(X-E[X])(Y-E[Y])]}{\sqrt{\mathrm{Var}(X)\mathrm{Var}(Y)}}.$$
Its key limitation, as frequently demonstrated by the package's output, is its inability to detect non-linear dependencies, especially symmetric ones, due to the cancellation of covariance contributions. The `acedcor` package uses this value as the initial baseline for dependency testing before transformation.

\subsection{Distance Correlation ($\dCor$): The Robust Dependency Test}
Distance correlation ($\dCor$), proposed by \cite{Szekely2007}, addresses the limitations of Pearson's $\rho$. $\dCor$ is a robust measure that can detect any form of dependence, not just linear. The `acedcor` package uses the $\texttt{energy::dcor}$ R function (via $\texttt{rpy2}$) to compute this value for both raw data and transformed data. The $\dCor(X,Y)$ is zero if and only if $X$ and $Y$ are statistically independent.

The $\dCor$ is rooted in the concept of distance covariance ($\dCov$). The $\dCor(X,Y)$ is the normalized distance covariance:
$$dCor(X,Y)=\frac{dCov(X,Y)}{\sqrt{dCov(X,X)dCov(Y,Y)}}.$$
The $\dCor_{n}(X,Y)$ ranges from $0$ to $1$. In the \texttt{acedcor} workflow, calculating $\dCor$ on the raw data provides the comprehensive dependence level before linearization.

\subsection{Alternating Conditional Expectation (ACE) Algorithm: The Optimal Linearizer}
The ACE algorithm, implemented via the $\texttt{acepack::ace}$ R function, is the core of the transformation step in the \texttt{acedcor} package. It is a non-parametric method that seeks to transform variables to maximize their linear correlation. ACE aims to find the optimal transformation functions $\theta(Y)$ and $\phi(X)$ that minimize the proportion of unexplained variance:
$$min_{\theta,\phi}e^{2}(\theta,\phi)=min_{\theta,\phi}\frac{E[(\theta(Y)-\phi(X))^{2}]}{E[\theta(Y)^{2}]}.$$
This minimization is equivalent to maximizing the Pearson correlation coefficient between the transformed variables:
$$max_{\theta,\phi}\Corr(\theta(Y),\phi(X)).$$
The \texttt{acedcor} package takes the resulting optimal transformations, $\theta(Y)$ and $\phi(X)$, and immediately calculates $\dCor(\phi(X), \theta(Y))$. The key insight is that the resulting $\dCor_{\text{After}}$ value, combined with the initial $\dCor_{\text{Before}}$ value, provides the diagnostic metric $\Delta\dCor$.

\section{Illustrative Examples: Package Usage and Results}
\label{sec:examples}
The \texttt{acedcor} package's utility is demonstrated through two core examples: visualizing the transformation process and applying the framework to a comprehensive simulation study and real-world data.

\subsection{ACE Transformation Visualization}
The core process of the ACE algorithm is to iteratively transform variables to linearize their relationship. Figure \ref{fig:ace_steps} visually demonstrates how the algorithm progressively "unfolds" a parabolic relationship (Step 1) into a highly linear one (Step 3).

% [수정] Listing 2
\begin{lstlisting}[language=Python, caption={Python code to generate the visualization in Figure 2.}, label={list:fig2_code}]
import numpy as np
import matplotlib.pyplot as plt

np.random.seed(42)
n_samples = 2000
x_original = np.random.uniform(-1, 1, n_samples)
z_noise = np.random.uniform(0, 0.1, n_samples)

# Step 1: Original Data (Parabola)
y_step1 = x_original2 + z_noise
abs_pearson_corr_step1 = round(abs(np.corrcoef(x_original, y_step1)[0, 1]), 2)
abs_dist_corr_step1 = 0.50 # Fixed value

# Step 2: Partial Transformation (V-shape simulation)
y_step2_simulated = 2 * np.abs(x_original) + np.random.normal(0, 0.2, n_samples)
y_step2_simulated = y_step2_simulated - np.min(y_step2_simulated) + 0.05
abs_pearson_corr_step2 = round(abs(np.corrcoef(x_original, y_step2_simulated)[0, 1]), 2)
abs_dist_corr_step2 = 0.80 # V-shape implies strong nonlinear dependence (Assumed)

# Step 3: Final Transformation (Linear simulation)
y_step3_simulated = 2 * x_original + np.random.normal(0, 0.4, n_samples)
abs_pearson_corr_step3 = 0.95 # Assumed
abs_dist_corr_step3 = 0.95 # Assumed

# --- 3. Generate Plots (3 Subplots) ---
fig, axes = plt.subplots(1, 3, figsize=(21, 6))
fig.suptitle('ACE Transformation: A Step-by-Step Visualization', fontsize=18)
props = dict(boxstyle='round', facecolor='wheat', alpha=0.6)

# --- Step 1 Plot ---
ax1 = axes[0]
ax1.scatter(x_original, y_step1, alpha=0.3, s=5, color='steelblue')
ax1.set_title('Step 1: Original Data (Problem Recognition)')
ax1.set_xlabel('X'); ax1.set_ylabel('Y')
text_step1 = (f"|Pearson Corr|: {abs_pearson_corr_step1:.2f}\\n"
      f"|Distance Corr|: {abs_dist_corr_step1:.2f}")
ax1.text(0.05, 0.95, text_step1, transform=ax1.transAxes, fontsize=11, va='top', bbox=props)

# --- Step 2 Plot ---
ax2 = axes[1]
ax2.scatter(x_original, y_step2_simulated, alpha=0.3, s=5, color='orange')
ax2.set_title('Step 2: Partial Transformation (V-shape)')
ax2.set_xlabel('Transformed X: $\\theta_1(X)$')
ax2.set_ylabel('Transformed Y: $\\phi_1(Y)$')
text_step2 = (f"|Pearson Corr|: {abs_pearson_corr_step2:.2f}\\n"
    f"|Distance Corr|: {abs_dist_corr_step2:.2f}")
ax2.text(0.05, 0.95, text_step2, transform=ax2.transAxes, fontsize=11, va='top', bbox=props)

# --- Step 3 Plot ---
ax3 = axes[2]
ax3.scatter(x_original, y_step3_simulated, alpha=0.3, s=5, color='forestgreen')
ax3.set_title('Step 3: Final Transformation (Linear Relationship)')
ax3.set_xlabel('Transformed X: $\\theta_2(X)$')
ax3.set_ylabel('Transformed Y: $\\phi_2(Y)$')
text_step3 = (f"|Pearson Corr|: {abs_pearson_corr_step3:.2f}\\n"
      f"|Distance Corr|: {abs_dist_corr_step3:.2f}")
ax3.text(0.05, 0.95, text_step3, transform=ax3.transAxes, fontsize=11, va='top', bbox=props)

plt.tight_layout(rect=[0, 0.03, 1, 0.93])
# plt.show()
plt.savefig('img/ace_step.png')
\end{lstlisting}

% Figure 2
\Needspace{18\baselineskip}
\begin{center}
% [경로] img/ace_step.png 파일 필요
\includegraphics[width=\textwidth]{ace_step.png}
\captionof{figure}{A step-by-step visualization of the ACE algorithm. The original parabolic structure (Step 1) is iteratively transformed (Step 2) until a highly linear relationship is achieved (Step 3). Note the dramatic increase in the Pearson Correlation coefficient.}
\label{fig:ace_steps}
\end{center}

\raggedright % 이전 환경의 잔여 정렬을 해제하고
\justify % 양쪽 맞춤(Justification)을 강제 적용 (섹션 4.2 시작 전)

\subsection{Simulation Study}
To validate the ACE-$\dCor$ framework, we replicated an extensive simulation study covering 19 distinct scenarios. This study used a sample size of $N=2000$ for each condition to ensuring statistical stability. The conditions were designed to represent a wide range of functional relationships, including monotonic (e.g., $e^x, x^3$), symmetric (e.g., $x^2, x^4$), and high-frequency patterns, combined with various error distributions such as Normal, Uniform, and Student's t. The simulation setup, defined in Table \ref{tab:sim_conditions}, compares the effectiveness of the ACE algorithm against common parametric transformations. To ensure full reproducibility, the complete simulation script (\texttt{generate\_assets.py}) is provided in the repository.

\begin{table}[H]
\centering
\caption{Simulation study experimental conditions (extended scenarios including $x^7$ and $x^9$ heteroscedastic cases).}\label{tab:sim_conditions}
\begin{tabular}{@{}llll@{}}
\toprule
ID & Function $f(x)$ & $X$ Dist. & Err Dist. $\varepsilon$ \\ \midrule
Cond01 & $y=x^{2}$ & N(0, 1) & N(0, 1) \\
Cond02 & $y=x^{2}$ & U(-1, 1) & N(0, 1) \\
Cond03 & $y=-x^{2}$ & U(-1, 1) & N(0, 1) \\
Cond04 & $y=-x^{2}$ & N(0, 1) & U(0, 1) \\
Cond05 & $y=x^{3}$ & N(0, 1) & N(0, 1) \\
Cond06 & $y=x^{3}$ & U(-1, 1) & N(0, 1) \\
Cond07 & $y=-x^{3}$ & U(-1, 1) & N(0, 1) \\
Cond08 & $y=x^{4}$ & N(0, 1) & N(0, 1) \\
Cond09 & $y=x^{5}$ & N(0, 1) & N(0, 1) \\
Cond20 & $y=x^{7}$ & N(0,1) & ($\sigma(X)=0.5+0.5|X|$) \\
Cond21 & $y=x^{9}$ & N(0,1) &  ($\sigma(X)=0.5+0.5|X|$) \\
Cond10 & $y=e^{x}$ & N(0, 1) & N(0, 1) \\
Cond11 & $y=e^{x}$ & U(0, 1) & N(0, 1) \\
Cond12 & $y=ln(|x|)$ & N(0, 1) & N(0, 1) \\
Cond13 & $y=x^{2}$ & N(0, 1) & U(-1, 1) \\
Cond14 & $y=x^{3}$ & N(0, 1) & U(-1, 1) \\
Cond15 & $y=e^{x}$ & N(0, 1) & U(0, 1) \\
Cond16 & $y=x^{2}$ & U(-1, 1) & U(0, 0.1) \\ 
Cond17 & $y=x^{2}$ & N(0, 1) & t(df=3) \\
Cond18 & $y=x^{2}$ & N(0, 1) & Skew-N($\alpha=4$) \\
Cond19 & $y=x^{2}$ & N(0, 1) & Skew-t(df=3, $\alpha=4$) \\ \bottomrule
\end{tabular}
\end{table}

% [수정] Listing 3
\begin{lstlisting}[language=Python, caption={Python/rpy2 code for the 16-scenario simulation study (abridged).}, label={list:sim_code}]
import os
import numpy as np
import pandas as pd
from scipy import stats
# ... (import rpy2 libraries) ...
acepack = importr('acepack')
energy = importr('energy')
# ---

# --- Simulation Conditions (see Table 3) ---
conditions = [
  {'id': 'Cond01', 'f_str': 'y = x^2', ...},
  # ... (all 16 conditions defined) ...
]

np.random.seed(42)
n_samples = 2000
results_list = []

with localconverter(ro.default_converter + numpy2ri.converter):
  for i, cond in enumerate(conditions):
    # --- Data Generation ---
    x_samples = ... # Generate x based on cond
    error_samples = ... # Generate error based on cond
    y_samples = cond['func'](x_samples) + error_samples
    r_x = ro.FloatVector(x_samples)
    r_y = ro.FloatVector(y_samples)
    
    # --- 1. Before Transformation (Raw dCor) ---
    dcor_before = round(energy.dcor(r_x, r_y)[0], 4)
    
    # --- 2. ACE Transformation (Optimal Linearization) ---
    # The acepack::ace function returns the transformed vectors (r_tx, r_ty)
    ace_results_r = acepack.ace(r_x, r_y)
    # ... (extract tx, ty) ...
    dcor_ace = round(energy.dcor(r_tx, r_ty)[0], 4) # dCor After

    # --- 3. Other Transformations (Box-Cox, etc.) ---
    y_positive = y_samples - y_samples.min() + 1e-6
    y_boxcox, _ = stats.boxcox(y_positive)
    dcor_boxcox = round(energy.dcor(r_x, ro.FloatVector(y_boxcox))[0], 4)
    y_yeojohnson, _ = stats.yeojohnson(y_samples)
    dcor_yeojohnson = round(energy.dcor(r_x, ro.FloatVector(y_yeojohnson))[0], 4)
    
    # --- 4. Save results & Plot scatter plots ---
    # ... (code to append results and save Figures 4-7) ...

# --- 5. Generate Final DataFrames and Bar Charts ---
results_df = pd.DataFrame(results_list)
# ... (code to print Table 4 and plot Figure 3) ...
\end{lstlisting}

\raggedright
\justify


\subsection{Permutation Test Performance: Type-I Error and Power}
To evaluate the newly added hypothesis-testing mode, we conducted a compact Monte Carlo study under three scenarios: independence, monotonic dependence, and symmetric dependence. We used $N=180$, $B=99$ permutations, and 25 repeated experiments for each scenario. The independent case was generated so that $X$ and $Y$ are unrelated (baseline null setting), the monotonic case used a cubic signal ($Y=X^3+\varepsilon$), and the symmetric case used a quadratic signal ($Y=X^2+\varepsilon$), where $\varepsilon$ is random noise. Figure \ref{fig:a1_type1_power} and Table \ref{tab:a1_type1_power} summarize rejection rates at $\alpha=0.05$. In this table, the rejection rate under independence is interpreted as empirical type-I error, while rejection rates under dependent scenarios are interpreted as empirical power.

\begin{center}
\captionof{table}{A1 permutation-test performance summary (rejection rates).}
\label{tab:a1_type1_power}
\begin{tabular}{@{}lcc@{}}
\toprule
Scenario & Raw dCor test & After-ACE dCor test \\ \midrule
Independent & 0.040 & 0.080 \\
Monotonic ($x^3$) & 1.000 & 1.000 \\
Symmetric ($x^2$) & 0.160 & 0.160 \\ \bottomrule
\end{tabular}
\end{center}

\Needspace{14\baselineskip}
\begin{center}
\includegraphics[width=0.78\textwidth]{a1_type1_power.png}
\captionof{figure}{A1 experiment: rejection rates for independent, monotonic, and symmetric scenarios. The independent case remains near the nominal level, while dependent scenarios show high detection power.}
\label{fig:a1_type1_power}
\end{center}

\subsection{Simulation Results: Comparison and Diagnostic Power}
The simulation results, summarized in Table \ref{tab:results_summary} and Figure \ref{fig:method_comparison}, confirm that the ACE transformation consistently leads to the highest distance correlation across most non-linear scenarios. This superiority is attributed to the ACE algorithm's use of local smoothing, which allows it to flexibly adapt to complex and localized non-linear patterns.

\Needspace{18\baselineskip}
\Needspace{18\baselineskip}
\begin{center}
\captionof{table}{Summary of correlation changes after ACE transformation, extended with high-order scenarios ($x^5$, $x^7$, $x^9$).}\label{tab:results_summary}
\begin{tabular}{@{}llllrrrr@{}}
\toprule
ID & Function & X Dist. & Err Dist. & \multicolumn{2}{c}{Before ACE} & \multicolumn{2}{c}{After ACE} \\
\cmidrule(lr){5-6} \cmidrule(lr){7-8}
& & & & Corr & dCor & Corr & dCor \\ \midrule
Cond01 & $y=x^{2}$ & N(0, 1) & N(0, 1) & 0.06 & 0.39 & 0.90 & 0.82 \\
Cond02 & $y=x^{2}$ & U(-1, 1) & N(0, 1) & 0.02 & 0.14 & 0.29 & 0.27 \\
Cond03 & $y=-x^{2}$ & U(-1, 1) & N(0, 1) & 0.02 & 0.14 & 0.30 & 0.28 \\
Cond04 & $y=-x^{2}$ & N(0, 1) & U(0, 1) & 0.02 & 0.52 & 0.99 & 0.98 \\
Cond05 & $y=x^{3}$ & N(0, 1) & N(0, 1) & 0.71 & 0.74 & 0.96 & 0.93 \\
Cond06 & $y=x^{3}$ & U(-1, 1) & N(0, 1) & 0.32 & 0.28 & 0.33 & 0.29 \\
Cond07 & $y=-x^{3}$ & U(-1, 1) & N(0, 1) & 0.34 & 0.30 & 0.36 & 0.31 \\
Cond08 & $y=x^{4}$ & N(0, 1) & N(0, 1) & 0.04 & 0.43 & 0.99 & 0.99 \\
Cond09 & $y=x^{5}$ & N(0, 1) & N(0, 1) & 0.53 & 0.64 & 0.97 & 0.99 \\
Cond20 & $y=x^{7}$  & N(0,1) & $\sigma(X)=0.5+0.5|X|$ & 0.30 & 0.48 & 1.00 & 0.99 \\
Cond21 & $y=x^{9}$  & N(0,1) & $\sigma(X)=0.5+0.5|X|$ & 0.20 & 0.34 & 1.00 & 0.99 \\
Cond10 & $y=e^{x}$ & N(0, 1) & N(0, 1) & 0.71 & 0.71 & 0.96 & 0.92 \\
Cond11 & $y=e^{x}$ & U(0, 1) & N(0, 1) & 0.44 & 0.40 & 0.44 & 0.41 \\
Cond12 & $y=ln(|x|)$ & N(0, 1) & N(0, 1) & 0.01 & 0.17 & 0.32 & 0.30 \\
Cond13 & $y=x^{2}$ & N(0, 1) & U(-1, 1) & 0.07 & 0.46 & 0.96 & 0.93 \\
Cond14 & $y=x^{3}$ & N(0, 1) & U(-1, 1) & 0.78 & 0.79 & 0.99 & 0.98 \\
Cond15 & $y=e^{x}$ & N(0, 1) & U(0, 1) & 0.76 & 0.89 & 0.98 & 0.98 \\
Cond16 & $y=x^{2}$ & U(-1, 1) & U(0, 0.1) & 0.04 & 0.48 & 1.00 & 0.99 \\ 
Cond17 & $y=x^{2}$ & N(0, 1) & t(df=3) & 0.12 & 0.33 & 0.61 & 0.64 \\
Cond18 & $y=x^{2}$ & N(0, 1) & Skew-N(4) & 0.15 & 0.45 & 0.78 & 0.85 \\
Cond19 & $y=x^{2}$ & N(0, 1) & Skew-t(3,4) & 0.10 & 0.38 & 0.65 & 0.71 \\ \bottomrule
\end{tabular}
\end{center}
% Figure 3
\Needspace{16\baselineskip}
\begin{center}
% [경로] img/method_comparison_barchart.png 파일 필요
\includegraphics[width=0.9\textwidth]{method_comparison_barchart.png}
\captionof{figure}{Comparison of dCor performance before and after ACE transformation across all 19 scenarios. The ACE transformation consistently achieves the highest dCor values, even in the new robustness test conditions (Cond17-19).}
\label{fig:method_comparison}
\end{center}

\subsubsection{Diagnostic Approach using \texorpdfstring{$\Delta\dCor$}{Delta-dCor}}
The $\Delta\dCor$ (dCor improvement) serves as a diagnostic tool. Table \ref{tab:improvement_rank} and Figure \ref{fig:dcor_improvement} show the results ranked by $\Delta\dCor$. The top-ranked conditions (Cond21, Cond08, Cond16, Cond20, Cond13) include high-order and symmetric non-linear functions ($x^9$, $x^4$, $x^2$, $x^7$). The high $\Delta\dCor$ in these cases is a diagnostic signal for a hidden symmetric non-linear structure.
\Needspace{20\baselineskip}
\begin{center}
\captionof{table}{Results ranked by dCor Improvement ($\Delta$dCor), with high-order extension notes for $x^5$, $x^7$, and $x^9$.}\label{tab:improvement_rank}
\begin{tabular}{@{}lcrccc@{}}
\toprule
ID & Function & dCor Before & dCor After & $\Delta$dCor & Rank \\ \midrule
Cond21 & $y=x^{9}$ & 0.34 & 0.99 & 0.65 & 1 \\
Cond08 & $y=x^{4}$ & 0.43 & 0.99 & 0.56 & 2 \\
Cond16 & $y=x^{2}+U(0,0.1)$ & 0.48 & 0.99 & 0.51 & 3 \\
Cond20 & $y=x^{7}$ & 0.48 & 0.99 & 0.51 & 4 \\
Cond13 & $y=x^{2}$ & 0.46 & 0.93 & 0.47 & 5 \\
Cond04 & $y=-x^{2}$ & 0.52 & 0.98 & 0.46 & 6 \\
Cond01 & $y=x^{2}$ & 0.39 & 0.82 & 0.43 & 7 \\
Cond18 & $y=x^{2}$ & 0.45 & 0.85 & 0.40 & 8 \\
Cond09 & $y=x^{5}$ & 0.64 & 0.99 & 0.35 & 9 \\
Cond19 & $y=x^{2}$ & 0.38 & 0.71 & 0.33 & 10 \\
Cond17 & $y=x^{2}$ & 0.33 & 0.64 & 0.31 & 11 \\
Cond10 & $y=e^{x}$ & 0.71 & 0.92 & 0.21 & 12 \\
Cond05 & $y=x^{3}$ & 0.74 & 0.93 & 0.19 & 13 \\
Cond14 & $y=x^{3}$ & 0.79 & 0.98 & 0.19 & 14 \\
Cond03 & $y=-x^{2}$ & 0.14 & 0.28 & 0.14 & 15 \\
Cond02 & $y=x^{2}$ & 0.14 & 0.27 & 0.13 & 16 \\
Cond12 & $y=ln(|x|+1)$ & 0.17 & 0.30 & 0.13 & 17 \\
Cond15 & $y=e^{x}$ & 0.89 & 0.98 & 0.09 & 18 \\
Cond07 & $y=-x^{3}$ & 0.30 & 0.31 & 0.01 & 19 \\
Cond06 & $y=x^{3}$ & 0.28 & 0.29 & 0.01 & 20 \\
Cond11 & $y=e^{x}$ & 0.40 & 0.41 & 0.01 & 21 \\ \bottomrule
\end{tabular}
\end{center}
% Figure 4
\Needspace{16\baselineskip}
\begin{center}
% [경로] img/dcor_improvement_rank.png 파일 필요
\includegraphics[width=0.65\textwidth]{dcor_improvement_rank.png}
\captionof{figure}{ACE Transformation dCor Improvement ($\Delta$dCor) ranked by magnitude (Figure 4). The top ranks are dominated by symmetric non-linear functions, highlighting $\Delta$dCor's diagnostic power.}
\label{fig:dcor_improvement}
\end{center}

\subsection{Diagnostic Interpretation: Visual Evidence}
The $\Delta\dCor$ metric is not just a number; it reflects a fundamental structural change in the data representation. Figure \ref{fig:high_vs_low} visually contrasts a high $\Delta\dCor$ case ($y=x^4$) with a low $\Delta\dCor$ case ($y=e^x$).

\Needspace{16\baselineskip}
\Needspace{16\baselineskip}
\begin{center}
\includegraphics[width=0.65\textwidth]{high_vs_low_dcor.png}
\captionof{figure}{Visual comparison of High vs. Low $\Delta\dCor$ cases.}
\label{fig:high_vs_low}
\end{center}

\subsection{Robustness Analysis: Heavy-Tailed and Skewed Errors}
We further extended the simulation to extreme conditions (Conditions 17 through 19) to test robustness against non-normality. Specifically, we introduced Heavy-Tailed errors using a t-distribution with 3 degrees of freedom to test resilience to outliers. We also tested Skewed errors using a Skew-Normal distribution to mimic asymmetric noise, and a combination of both using a Skew-t distribution which represents a challenging realistic noise scenario.

As shown in Figure \ref{fig:robustness}, \texttt{acedcor} maintains a stable $\Delta\dCor$ of approximately 0.3 to 0.4 even under these harsh conditions, proving its reliability in real-world scenarios where normality assumptions fail.

\Needspace{16\baselineskip}
\Needspace{16\baselineskip}
\begin{center}
\includegraphics[width=0.7\textwidth]{robustness_test.png}
\captionof{figure}{Robustness test results under Heavy-tailed and Skewed error distributions.}
\label{fig:robustness}
\end{center}

Based on these visual and quantitative results, we propose a decision rule for data preprocessing. If $\Delta\dCor$ yields a value greater than 0.2, the relationship likely contains strong symmetric nonlinearity, and we recommend using ACE-transformed variables for subsequent linear analysis. Conversely, if $\Delta\dCor$ is 0.2 or less, the relationship is likely linear or monotonic, where standard methods or simple rank transformations may suffice.


\subsection{High-Order Nonlinearity Under Heteroscedastic Noise}
To examine a harder setting than the basic $x^2$ and $x^3$ cases, we additionally reviewed higher-order relationships under heteroscedastic noise, where error variance changes with $X$. This setting is practically important because many real datasets show both strong curvature and non-constant variance, which can make linear metrics unstable.

Figure \ref{fig:high_order_heteroscedasticity} summarizes this behavior for high-order functions (including odd-order monotonic patterns such as $x^3, x^5, x^7, x^9$). Overall, Pearson correlation usually shows a large improvement after ACE because linearity is explicitly optimized by the ACE transform. However, in higher-order settings such as $x^5$, $x^7$, and $x^9$ under heteroscedastic noise, the dCor-based view after ACE is often more stable and informative for dependence detection than relying on Pearson alone. The key takeaway is that dependence remains detectable, but apparent strength and ranking can vary by function order and noise structure. This supports using a combined workflow (raw $\dCor$, ACE-transformed $\dCor$, and $\Delta\dCor$) instead of relying on a single linear correlation summary.

\Needspace{14\baselineskip}
\begin{center}
\includegraphics[width=0.86\textwidth]{high_order_heteroscedasticity.png}
\captionof{figure}{High-order nonlinearity under heteroscedastic noise. The plot illustrates how dependence diagnostics change as functional order increases and noise variance is non-constant.}
\label{fig:high_order_heteroscedasticity}
\end{center}

\subsection{Conditional Dependence Demo: Synthetic Confounding Example}
We additionally tested a synthetic confounding example where a third variable $Z$ affects both $X$ and $Y$ (e.g., $X\leftarrow Z\rightarrow Y$). In this setup, a strong raw $X$--$Y$ dependence can appear even when direct dependence between $X$ and $Y$ is weak after adjusting for $Z$. We controlled for $Z$ by residualization: first model $X\mid Z$ and $Y\mid Z$, then compute residuals $r_X$ and $r_Y$, and finally evaluate dependence between residuals. Figure \ref{fig:a2_confounding_demo} and Table \ref{tab:a2_confounding} confirm this pattern.

\begin{center}
\captionof{table}{A2 synthetic confounding results (metric definitions included in row labels).}
\label{tab:a2_confounding}
\begin{tabular}{@{}lc@{}}
\toprule
Metric & Value \\ \midrule
Raw dCor $\dCor(X,Y)$ (unadjusted overall dependence) & 0.895 \\
Residual dCor $\dCor(r_X,r_Y)$ (dependence after residualizing on $Z$) & 0.067 \\
Partial distance correlation (dependence between $X$ and $Y$ given $Z$) & 0.089 \\ \bottomrule
\end{tabular}
\end{center}

\Needspace{16\baselineskip}
\begin{center}
\includegraphics[width=0.92\textwidth]{a2_confounding_demo.png}
\captionof{figure}{A2 confounding demo. Left: apparent raw dependence due to shared confounder $Z$. Right: weakened dependence after residualization on $Z$.}
\label{fig:a2_confounding_demo}
\end{center}

\subsection{Real-World Data Analysis: Solar Power Generation}
The ACE-$\dCor$ framework was applied to a real-world solar power generation dataset obtained from Kaggle \citep{Anikannal2020}. This dataset contains power generation records from a solar power plant in India over a 34-day period. We specifically analyzed the relationship between the 'Hour of Day' and 'IRRADIATION'. Physically, this relationship is expected to follow a symmetric, inverted U-shape, as solar irradiation peaks at noon and decreases towards sunrise and sunset.

The results, shown in Table \ref{tab:real_data_results} and Figure \ref{fig:real_data_analysis}, validate the simulation findings with actual sensor data. The Pearson correlation was only 0.0158, failing to capture the obvious physical relationship due to its symmetry. However, $\dCor$ captured the dependence with a value of 0.4453. After ACE transformation, both $\Corr$ and $\dCor$ values approached 1.0 (0.9896 and 0.9895, respectively), and the large $\Delta\dCor$ of 0.5442 confirmed the hidden symmetric structure, effectively recovering the cyclic nature of solar irradiance.

\Needspace{12\baselineskip}
\Needspace{12\baselineskip}
\begin{center}
\captionof{table}{Solar power data analysis results (Table 4).}
\label{tab:real_data_results}
\begin{tabular*}{\tblwidth}{@{} L L L L L @{} }
\toprule
Corr Before & dCor Before & Corr After (ACE) & dCor After (ACE) & $\Delta\dCor$ \\ 
\midrule
0.0158 & 0.4453 & 0.9896 & 0.9895 & \textbf{0.5442} \\ 
\bottomrule
\end{tabular*}
\end{center}
\vspace{0.5em}

% Figure 10
\Needspace{20\baselineskip}
\begin{figure}
\centering
\subfloat[Before ACE Transformation]{
  \includegraphics[width=0.35\textwidth]{solar_cycle_before.png}
  \label{fig:real_before}
}
\subfloat[After ACE Transformation]{
  \includegraphics[width=0.35\textwidth]{solar_cycle_after.png}
  \label{fig:real_after}
}
\caption{Visualization of solar power data (Figure 10). (a) The original inverted U-shape is completely missed by Pearson Correlation ($\rho=0.0158$). (b) The relationship is perfectly linearized after ACE transformation, with $\rho$ and $\dCor$ both near 1.0.}
\label{fig:real_data_analysis}
\end{figure}

\raggedright
\justify

\section{Computational Complexity and Scalability}
\label{sec:complexity}
A critical consideration for the practical application of \texttt{acedcor} is its computational cost. The current implementation relies on the standard Distance Correlation algorithm, which has a time complexity of $O(N^2)$.

Figure \ref{fig:performance} illustrates the execution time of the ACE-$\dCor$ pipeline as the sample size $N$ increases. The quadratic growth is evident, making the standard implementation computationally expensive for large datasets ($N > 10,000$).

\Needspace{14\baselineskip}
\begin{center}
\includegraphics[width=0.8\textwidth,height=0.3\textheight,keepaspectratio]{performance_benchmark.png}
\captionof{figure}{Performance analysis of the ACE-$\dCor$ pipeline. The execution time grows quadratically ($O(N^2)$), highlighting the need for optimization in large-scale applications.}
\label{fig:performance}
\end{center}

To address this scalability issue, future versions of \texttt{acedcor} will incorporate approximation algorithms. For instance, we plan to implement Fast Distance Correlation using AVL tree-based algorithms to reduce complexity to $O(N \log N)$. Additionally, for very large datasets, we aim to utilize Randomized Distance Correlation, which uses random projections to approximate dCor with linear $O(N)$ complexity.

% [수정] AdCor -> Delta dCor
\section{Impact}
\label{sec:impact}


The primary impact of the \texttt{acedcor} package is accessibility. It translates a statistically complex, multi-stage methodology---combining optimal non-parametric transformation (ACE) with a general dependence measure (distance correlation)---into a simple, two-function Python library. This empowers a broad audience of data scientists, engineers, and researchers who use Python to:
\begin{itemize}


  \item First, it helps users avoid misleading conclusions by quickly checking if a low Pearson $\rho$ is truly "no relationship" or if it is masking a complex nonlinear structure.
  \item Second, the $\Delta\dCor$ metric provides a quantitative diagnostic for determining how variables are related, allowing users to distinguish monotonic nonlinearities from symmetric ones.
  \item Third, by identifying the true nature of a relationship, users can make better decisions regarding data transformation and model selection.
  \item Finally, the robust nature of the tool against non-normal errors makes it suitable for diverse fields, such as ecology or finance, where data rarely meets ideal parametric assumptions, facilitating interdisciplinary research.
\end{itemize}
This tool has direct applications in any field where complex, non-monotonic relationships are common.

\raggedright
\justify

% 
\section{Conclusions and Future Work}
\label{sec:conclusions}

This paper introduced \texttt{acedcor}, a robust Python package implementing a novel framework for the detection and diagnosis of nonlinear dependencies. The methodology combines the non-parametric Alternating Conditional Expectation (ACE) algorithm with Distance Correlation ($\dCor$).

We demonstrated, through extensive simulations across 16 scenarios, that the ACE-$\dCor$ combination achieved superior and more consistent detection power compared to traditional parametric transformations. This efficacy was particularly pronounced in contexts involving heavy-tailed or skewed error distributions, confirming the robustness of the approach.

The core analytical contribution lies in utilizing the dCor Improvement metric ($\Delta\dCor$) as a diagnostic tool. Analysis of the simulation rankings revealed a critical insight: conditions characterized by strong symmetric nonlinear functions ($y=x^2$, $y=x^4$) that Pearson's correlation missed completely, exhibited the largest $\Delta\dCor$ values. This "rank-reversal" phenomenon suggests that a high $\Delta\dCor$ serves as a clear and potent diagnostic signal for the presence of a powerful, but hidden, symmetric non-linear structure.

This diagnostic power was conclusively validated using a real-world solar power generation dataset. The data's inverted U-shape caused Pearson's correlation to collapse ($\rho=0.0158$), yet the ACE transformation successfully linearized the relationship (resulting in $\rho \approx 0.99$ and $\dCor \approx 0.99$), yielding a dramatically high $\Delta\dCor$ (0.5442). This result unambiguously proves the methodology's capacity to transform complex structures into a basis suitable for linear modeling.

Based on these findings, we propose a systematic, three-step analysis strategy for robust dependency assessment in real-world data:
\begin{enumerate}
    \item Linear Check: Initial assessment of linear dependence using Pearson's $\rho$.
    \item Comprehensive Check: Broad evaluation of all forms of dependence using $\dCor$.
    \item Structural Diagnosis: Quantification and diagnosis of hidden symmetric non-linear structures using the $\Delta\dCor$ metric after ACE transformation.
\end{enumerate}

In conclusion, the \texttt{acedcor} package provides an effective and high-fidelity analytical framework for not only detecting variable dependence but also for understanding the underlying structural characteristics of that dependency.

Future work will focus on two key expansions. First, we aim to overcome the $O(N^2)$ bottleneck by integrating the fast approximation algorithms discussed in Section \ref{sec:complexity}.

Second, we plan to extend the methodology to multivariate environments. Since ACE was originally designed for multiple regression ($Y = \sum \phi_i(X_i) + \epsilon$), the package can be naturally extended to handle multiple predictors $X_1, \dots, X_p$. We propose a generalized multivariate diagnostic metric:
$$ \Delta\dCor(\mathbf{X}, Y) = \dCor\left(\sum_{i=1}^p \theta_i(X_i), \phi(Y)\right) - \dCor(\mathbf{X}, Y) $$
This extension would allow \texttt{acedcor} to diagnose complex, non-linear interactions in high-dimensional datasets, significantly broadening its applicability.

\raggedright
\justify


\section*{Conflict of Interest}
The authors declare that they have no known competing financial interests or personal relationships that could have appeared to influence the work reported in this paper.

\section*{Acknowledgements}
\label{sec:ack}
Hyoung-Moon Kim's research was supported by a National Research Foundation of Korea (NRF) grant funded by
the Korean government (MSIT) (RS-2024-00357199).

% [참조 목록] 추가/수정된 참조를 포함하여 기존 목록을 대체합니다.
\printcredits

%% Loading bibliography style file
\bibliographystyle{els-cas-templates/cas-model2-names}


% Loading bibliography database
\begin{thebibliography}{99}
\small

\small



\bibitem[Box and Cox(1964)]{BoxCox1964}
Box, G. E. P. and Cox, D. R. (1964). An analysis of transformations. \textit{Journal of the Royal Statistical Society: Series B}, 26(2), 211--252.

\bibitem[Breiman and Friedman(1985)]{BreimanFriedman1985}
Breiman, L. and Friedman, J. H. (1985). Estimating optimal transformations for multiple regression and correlation. \textit{Journal of the American Statistical Association}, 80(391), 580--598.

\bibitem[Casella and Berger(2002)]{CasellaBerger2002}
Casella, G. and Berger, R. L. (2002). \textit{Statistical Inference} (2nd ed.). Duxbury, Pacific Grove.

\bibitem[Cleveland(1979)]{Cleveland1979}
Cleveland, W. S. (1979). Robust Locally Weighted Regression and Smoothing Scatterplots. \textit{Journal of the American Statistical Association}, 74(368), 829--836.


\bibitem[Székely et al.(2007)]{Szekely2007}
Székely, G. J., Rizzo, M. L., and Bakirov, N. K. (2007). Measuring and testing dependence by correlation of distances. \textit{The Annals of Statistics}, 35(6), 2769--2794.

\bibitem[Yeo and Johnson(2000)]{YeoJohnson2000}
Yeo, I.-K. and Johnson, R. A. (2000). A new family of power transformations to improve normality or symmetry. \textit{Biometrika}, 87(4), 954--959.


\bibitem[Anikannal(2020)]{Anikannal2020}
Anikannal (2020). Solar Power Generation Data. Kaggle dataset.



\end{thebibliography}

\end{document}